% =========================================================================
% Extended Methods (Supplementary Information)
%
% Standalone document containing the technical detail offloaded from the
% main preprint.tex Methods section. The main text retains a one-line
% pointer at each offload site (``See Supplementary Methods'').
% =========================================================================

\documentclass[11pt,a4paper]{article}
\usepackage[utf8]{inputenc}
\usepackage[T1]{fontenc}
\usepackage{lmodern}
\usepackage[margin=2.5cm]{geometry}
\usepackage{microtype}
\usepackage{parskip}
\usepackage{amsmath,amssymb}
\usepackage{booktabs}
\usepackage{siunitx}
\DeclareSIUnit{\molar}{M}
\usepackage[colorlinks=true,linkcolor=blue,citecolor=blue,urlcolor=blue]{hyperref}
\usepackage{enumitem}
\usepackage[numbers,square,sort&compress]{natbib}

\title{Supplementary Methods\\
       \large Drug tolerance takes shape: persistent homology reveals
              cyclic cell-state plasticity in EGFR-mutant lung cancer}
\author{Hsieh-Ting Lin and Yueh-Hua Tu}
\date{}

\begin{document}
\maketitle

\section*{S1. Patient-atlas filtering details}

For GSE131907 \citep{kim2020} we filtered to cells with
\texttt{Sample\_Origin == `tLung'} (primary lung tumor) AND
\texttt{Cell\_type} containing `Epithelial', using the cell-type
annotations from \texttt{GSE131907\_Lung\_Cancer\_cell\_annotation.txt.gz}.
From the resulting 7{,}270 tumor epithelial cells, a uniform random
sample of 5{,}000 cells was drawn (\texttt{random\_state=42}), then
1{,}500 cells were subsampled further for tractable Vietoris--Rips
computation.

\section*{S2. Highly-variable gene (HVG) batch keys}

HVG selection used \texttt{batch\_key=`timepoint'} for GSE134839 and
GSE150949 (samples share a clonal PC9 background but differ in
drug-exposure timepoint, treated as a technical batch following
\citet{luecken2019}), \texttt{batch\_key=`pdx'} for GSE243562 (to
account for inter-PDX-model variability), and
\texttt{batch\_key=`Sample'} for the Kim atlas (per-patient
variability).

\section*{S3. PDX mouse-gene filter}

For PDX data, we filtered mouse genes using the pattern
\texttt{r'\^{}mm10.*|\^{}[a-z][a-z0-9]+\$'} applied to gene symbols
(matching either the Ensembl mouse prefix \texttt{mm10} or lowercase
gene symbols characteristic of MGI-style mouse annotations). Of
82{,}350 initial features, 24{,}062 mouse-like genes (29\%) were
removed, leaving 58{,}288 human-like features. We acknowledge this
symbol-based heuristic is inferior to alignment-based xeno-classification
(e.g.\ XenoCell, Xenomake) for publication-grade PDX analysis; our
downstream results depend on PCA structure and are insensitive to
modest filter errors, but this is flagged as a limitation in the main
text Discussion.

\section*{S4. Max $H_1$ as a hybrid statistic}

Max $H_1$ is a hybrid statistic: the existence of a nonzero $H_1$ bar
is a topological invariant (a nontrivial 1-cycle exists), but the
numerical value of its lifespan is measured in filtration-parameter
units (here, Euclidean distance in PCA space) and is therefore
metric-dependent. We refer to the statistic as ``topological'' because
it reports on the presence and robustness of a topological feature,
while acknowledging that the numerical scale depends on the metric
used. The PC-count sensitivity sweep (S5) quantifies the metric
dependence.

\section*{S5. PC-count sensitivity}

Max $H_1$ was computed on the D9 erlotinib timepoint ($n = 379$) and
D14 osimertinib timepoint ($n = 1{,}000$) at $n_{\text{PCs}} \in
\{10, 20, 30, 50\}$. Values were 1.96/2.08/1.77/1.03 for D9 and
3.12/5.56/3.78/4.40 for D14. The ordering drug-treated $>$ untreated
is preserved at every PC count, but absolute values vary within a
factor of 2. The $n_{\text{PCs}} = 30$ choice used throughout was
selected for consistency with prior scTDA work \citep{rizvi2017}.
See also Supplementary Table S5.

\section*{S6. EMT-cassette gene list (Mapper filter)}

EMT score (used as Mapper filter for the cell-line analysis) was
computed via \texttt{scanpy.tl.score\_genes} on a 33-gene MSigDB
Hallmark EMT cassette: \emph{VIM}, \emph{CDH2}, \emph{FN1},
\emph{SNAI1}, \emph{SNAI2}, \emph{ZEB1}, \emph{ZEB2}, \emph{TWIST1},
\emph{MMP2}, \emph{MMP9}, \emph{CDH1}, \emph{OCLN}, \emph{TJP1},
\emph{DSP}, \emph{KRT19}, \emph{KRT8}, \emph{SERPINE1}, \emph{TGFBI},
\emph{SPARC}, \emph{COL1A1}, \emph{COL3A1}, \emph{COL5A2},
\emph{ACTA2}, \emph{TAGLN}, \emph{MYL9}, \emph{CALD1}, \emph{LGALS1},
\emph{TPM1}, \emph{TPM2}, \emph{FSTL1}, \emph{ITGB1}, \emph{LAMC2},
\emph{LOXL2}. Of these, 28 of 33 genes were present in the GSE134839
normalized expression matrix and used for scoring.

\section*{S7. Harmony batch-correction parameters}

Harmony \citep{korsunsky2019} was applied to the top-50 PCA components
with default parameters: \texttt{theta = 2}, \texttt{lambda = 1},
\texttt{nclust = 100}, \texttt{epsilon\_cluster = 1e-5},
\texttt{epsilon\_harmony = 1e-4}, \texttt{max\_iter\_harmony = 10},
\texttt{max\_iter\_kmeans = 20}. Convergence was reached in 2
iterations for the osimertinib cohort. Batch keys: \texttt{timepoint}
for GSE150949, per-sample GSM identifier for GSE243562, per-patient
\texttt{Sample} identifier for GSE131907.

\section*{S8. Pipeline architecture}

The framework is organised as numbered modules, each accepting a
standard \texttt{AnnData} object and producing typed intermediate
outputs (persistence diagrams as \texttt{.npy}, Mapper graphs as
pickled dicts, enrichment results as \texttt{.csv}):

\begin{itemize}[leftmargin=1.25em]
  \item \texttt{01\_preprocess.py}: load GEO data, QC, normalise, HVG,
        regress, PCA, cell-cycle score.
  \item \texttt{02\_standard\_analysis.py}: UMAP, Leiden clustering
        \citep{traag2019} (multi-resolution), marker genes, EMT score,
        diffusion map, PAGA.
  \item \texttt{03\_topology.py}: per-group persistent homology, Mapper
        (3 filters), pairwise Wasserstein, permutation tests.
  \item \texttt{04\_cell\_cycle\_ablation.py}: remove cell-cycle genes,
        re-do PCA and PH, compute preservation ratio + permutation null.
  \item \texttt{05\_gene\_attribution.py}: Mapper gene connectivity,
        differential expression, pathway enrichment.
  \item \texttt{06\_validation.py}: apply the full pipeline to a
        validation dataset, compute cross-dataset Wasserstein.
  \item \texttt{07\_figures.py}: assemble publication figures from
        saved intermediate outputs.
  \item \texttt{08\_validation\_pdx.py}: PDX-specific loader and
        persistent homology for the ASCL1 PDX cohort.
  \item \texttt{09\_validation\_kim\_atlas.py}: Kim 2020 patient atlas
        baseline loading and $H_1$ computation.
  \item \texttt{10\_ablation\_permtest.py}: extended permutation tests
        on cell-cycle-ablated data.
  \item \texttt{11\_pdx\_gene\_attribution.py}: cross-system
        gene-overlap analysis between cell line and PDX.
  \item \texttt{12\_convert\_tables\_to\_parquet.py}: CSV $\to$ Parquet
        for large tables.
  \item \texttt{13\_harmony\_batch\_correction.py}: Harmony
        batch-correction sensitivity analysis.
  \item \texttt{14\_validation\_maynard.py}: Maynard 2020 neo-osi
        cohort loader.
  \item \texttt{15\_maynard\_full\_cohort.py}: full Maynard 14-patient
        EGFR-mutant cohort analysis.
\end{itemize}
Users can substitute their own datasets by providing an
\texttt{AnnData} file with a \texttt{timepoint} (or equivalent)
observation column.

\section*{S9. Unit tests}

The \texttt{sctda\_plasticity} package includes pytest-based unit tests
covering persistent homology on synthetic circle and cluster data,
Mapper graph construction with known topology, gene connectivity
ranking, and permutation test structure. All 8 tests pass on the
submitted version; see \texttt{tests/test\_topology.py}. Code is
lint-clean under \texttt{ruff check} with a 100-character line limit.

\section*{S10. Resampling stability of max $H_1$}

To characterise max $H_1$ variability under subsampling, we ran the
1{,}000-cell subsample five times with independent random seeds for
four cohorts (osi D14, osi D0, PDX YU-006 residual, PDX YU-006
untreated). Coefficients of variation (CV = SD/mean) were 8.9\%,
12.4\%, 12.8\%, and 23.5\% respectively. The \emph{ordering}
(D14 $>$ D0; residual $>$ untreated) was preserved in every resample.
Values in Table~\ref{tab:validation} of the main text are from seed
42; mean $\pm$ SD across the five seeds is in Supplementary
Table~S7.

\section*{S11. Alternative topological summaries}

Two alternative summaries of the persistence diagram on the
osimertinib time course replicate the monotonic trend. Total $H_1$
persistence (sum of finite bar lengths) grew $402 \to 568 \to 588$
across D0/D7/D14, while persistence entropy stayed at 6.7--7.0 ---
indicating the growth is driven by larger-scale features rather than
more numerous small ones. All three summaries across ten cohorts:
Supplementary Table~S6.

\section*{S12. Mapper parameter sweep}

A $4 \times 5$ Mapper-parameter sweep (\(n_{\text{cubes}} \in
\{10, 15, 20, 25, 30\}\); overlap \(\in \{0.2, 0.3, 0.4, 0.5\}\)) on
the cell-line dataset with the EMT filter produced 12--36 Mapper
nodes (Table~S3). Gene-connectivity rankings were near-identical
across all 20 combinations (mean Spearman \(\rho = 0.994\) vs.~the
reference, range 0.990--0.998; top-10 genes reproduced within rank
order 1--15; Supplementary Table~S8). The downstream EMT-enrichment
conclusion is therefore insensitive to the specific Mapper parameter
choice.

\section*{S13. Stringent S/G2M-score regression detail}

Because Tirosh ablation removes only a subset of the cell-cycle
program, we also applied the more demanding
\texttt{sc.pp.regress\_out(S\_score, G2M\_score)} on 1{,}000-cell
subsamples and recomputed max $H_1$. Across five cohorts (erlotinib
D0/D9/D11, PDX YU-006 untreated/residual), the post/pre ratio was
1.07--4.05 (Table~\ref{tab:stringent_cc}; mean 1.97, every value
$>1$), preserving --- and in several cases amplifying --- the loop
signal after the dominant cell-cycle variance axis is removed.
Erlotinib D11 ($n = 249$, ratio 4.05) is a case where regression
unmasks topological structure that PCA had obscured. The osimertinib
cohort contains zero Tirosh genes in its HVG set (filtered during
selection), so the Tirosh-specific contribution to its $H_1$ signal
is ruled out by absence; non-Tirosh cell-cycle genes (e.g.\
\emph{MKI67}, \emph{TOP2A}) do co-rank in PC1-Mapper, but the
converging discovery-cohort + PDX ablation evidence rules out
cell-cycle oscillation as the driver.

\section*{S14. Computational environment detail}

All analyses ran on macOS 15 (arm64) under Python 3.11 in a conda
environment specified in \texttt{envs/environment.yml}. Key package
versions: scanpy 1.11, anndata 0.12, ripser 0.6, persim 0.3,
kmapper 2.0, giotto-tda 0.6, scikit-learn 1.3, numpy 1.26,
pandas 2.x. The complete environment is reproducible via
\texttt{conda env create -f envs/environment.yml}. Random seed 42
is used for every stochastic operation (PCA, UMAP, Leiden
clustering, subsampling, permutation tests), yielding deterministic
outputs on the same platform up to numerical-library nondeterminism
(BLAS thread order, OpenMP scheduling). The framework has been
tested on both macOS (arm64) and Linux (x86\_64).

\bibliographystyle{unsrtnat}
\bibliography{../references}

\end{document}
